% E1.7c -- seleção de estrutura por limiarização: a saída (c) de
% docs/selecao-estrutura.md, decidida em D28.
%
% Consome o Corolário 4 de E1.6 (05-taxas.tex) e NADA MAIS: não usa
% irrepresentabilidade, não pede condição de desenho nova, e vale para o
% LASSO puro, que é o estimador base (D3). A saída (a), a propriedade
% oráculo da variante em grupos, é sondada em E1.7a e não é usada aqui.
%
% Numeração global dos resultados (docs/TAREFA.md, §5): este arquivo
% continua de Proposição 6, Lema 12, Teorema 2 e Corolário 7, e imprime o
% número global via \setcounter, como 04, 05 e 07.
%
% Conferência numérica: derivations/check/06-selecao-limiar.R (imprime OK).
% Notação: docs/notacao.md (congelada em E1.1). Macros: derivations/macros.tex.
\documentclass[11pt,a4paper]{article}
\usepackage[T1]{fontenc}
\usepackage[utf8]{inputenc}
\usepackage[brazil]{babel}
\usepackage[margin=2.6cm]{geometry}
\usepackage{enumitem}
\usepackage{booktabs}
\usepackage{xcolor}
\usepackage[colorlinks=true,linkcolor=blue!60!black,citecolor=blue!60!black]{hyperref}
% Macros do WAFC. Implementa docs/notacao.md (esboço; congela em E1.1).
% Único lugar onde uma macro é definida; os derivations/*.tex fazem % Macros do WAFC. Implementa docs/notacao.md (esboço; congela em E1.1).
% Único lugar onde uma macro é definida; os derivations/*.tex fazem % Macros do WAFC. Implementa docs/notacao.md (esboço; congela em E1.1).
% Único lugar onde uma macro é definida; os derivations/*.tex fazem \input{macros}.
\usepackage{amsmath,amssymb,amsthm}
\newcommand{\R}{\mathbb{R}}
\newcommand{\E}{\mathbb{E}}
\newcommand{\Prob}{\mathbb{P}}
\newcommand{\T}{^{\top}}
\newcommand{\norm}[1]{\left\lVert #1 \right\rVert}
\newcommand{\normn}[1]{\left\lVert #1 \right\rVert_{n}}
\newcommand{\normL}[1]{\left\lVert #1 \right\rVert_{L_2(P)}}
\newcommand{\Besov}[3]{B^{#1}_{#2,#3}}
\newcommand{\bX}{\mathbf{X}}
\newcommand{\bU}{\mathbf{U}}
\newcommand{\bZ}{\mathbf{Z}}
\newcommand{\btheta}{\boldsymbol{\theta}}
\newcommand{\thetastar}{\btheta^{*}}
\newcommand{\bc}{\mathbf{c}}
\newcommand{\PiJ}{\Pi_{J}}
\newcommand{\VJ}{V_{J}}
\newcommand{\NJ}{N_{J}}
\newtheorem{proposition}{Proposition}
\newtheorem{lemma}{Lemma}
\newtheorem{theorem}{Theorem}
\newtheorem{corollary}{Corollary}
\theoremstyle{definition}
\newtheorem{assumption}{Assumption}
\newtheorem{remark}{Remark}
.
\usepackage{amsmath,amssymb,amsthm}
\newcommand{\R}{\mathbb{R}}
\newcommand{\E}{\mathbb{E}}
\newcommand{\Prob}{\mathbb{P}}
\newcommand{\T}{^{\top}}
\newcommand{\norm}[1]{\left\lVert #1 \right\rVert}
\newcommand{\normn}[1]{\left\lVert #1 \right\rVert_{n}}
\newcommand{\normL}[1]{\left\lVert #1 \right\rVert_{L_2(P)}}
\newcommand{\Besov}[3]{B^{#1}_{#2,#3}}
\newcommand{\bX}{\mathbf{X}}
\newcommand{\bU}{\mathbf{U}}
\newcommand{\bZ}{\mathbf{Z}}
\newcommand{\btheta}{\boldsymbol{\theta}}
\newcommand{\thetastar}{\btheta^{*}}
\newcommand{\bc}{\mathbf{c}}
\newcommand{\PiJ}{\Pi_{J}}
\newcommand{\VJ}{V_{J}}
\newcommand{\NJ}{N_{J}}
\newtheorem{proposition}{Proposition}
\newtheorem{lemma}{Lemma}
\newtheorem{theorem}{Theorem}
\newtheorem{corollary}{Corollary}
\theoremstyle{definition}
\newtheorem{assumption}{Assumption}
\newtheorem{remark}{Remark}
.
\usepackage{amsmath,amssymb,amsthm}
\newcommand{\R}{\mathbb{R}}
\newcommand{\E}{\mathbb{E}}
\newcommand{\Prob}{\mathbb{P}}
\newcommand{\T}{^{\top}}
\newcommand{\norm}[1]{\left\lVert #1 \right\rVert}
\newcommand{\normn}[1]{\left\lVert #1 \right\rVert_{n}}
\newcommand{\normL}[1]{\left\lVert #1 \right\rVert_{L_2(P)}}
\newcommand{\Besov}[3]{B^{#1}_{#2,#3}}
\newcommand{\bX}{\mathbf{X}}
\newcommand{\bU}{\mathbf{U}}
\newcommand{\bZ}{\mathbf{Z}}
\newcommand{\btheta}{\boldsymbol{\theta}}
\newcommand{\thetastar}{\btheta^{*}}
\newcommand{\bc}{\mathbf{c}}
\newcommand{\PiJ}{\Pi_{J}}
\newcommand{\VJ}{V_{J}}
\newcommand{\NJ}{N_{J}}
\newtheorem{proposition}{Proposition}
\newtheorem{lemma}{Lemma}
\newtheorem{theorem}{Theorem}
\newtheorem{corollary}{Corollary}
\theoremstyle{definition}
\newtheorem{assumption}{Assumption}
\newtheorem{remark}{Remark}


% --- atalhos locais deste arquivo -----------------------------------------
\newcommand{\lmin}{\lambda_{\min}}
\newcommand{\Amat}{\mathbf{A}}
\newcommand{\Bmat}{\mathbf{B}}
\newcommand{\bb}{\mathbf{b}}
\newcommand{\bveps}{\boldsymbol{\varepsilon}}
\newcommand{\gtil}{\widetilde{\gamma}}
\newcommand{\Btil}{\widetilde{\Bmat}}
\newcommand{\bY}{\mathbf{Y}}
\newcommand{\bbeta}{\boldsymbol{\beta}}
\newcommand{\betahat}{\widehat{\bbeta}}
\newcommand{\smax}{\widehat{\sigma}_{\max}}
\newcommand{\eps}{\varepsilon}
\newcommand{\ghat}[2]{\widehat{g}_{#1#2}}
\newcommand{\Nhat}[2]{\widehat{N}_{#1#2}}
\newcommand{\Ntrue}[2]{N_{#1#2}}
\newcommand{\Shat}{\widehat{S}}
\newcommand{\sepmin}{\delta}
\newcommand{\rn}{\rho}

% A numeração global começa depois de Proposição 6, Lema 12, Teorema 2 e
% Corolário 7 (docs/TAREFA.md, §5; E1.8 usou o 6 e o 7).
\setcounter{proposition}{6}
\setcounter{lemma}{12}
\setcounter{theorem}{2}
\setcounter{corollary}{7}

% A hipótese deste arquivo é a única, e ganha a letra S para não colidir com
% a "Hipótese 1" de E1.6 (regime assintótico), que é citada ao lado dela.
\renewcommand{\theassumption}{S}

\title{E1.7c. Seleção de estrutura por limiarização}
\author{wafc-draft, \texttt{derivations/06-selecao-limiar.tex}}
\date{2026-09-20}

\begin{document}
\maketitle

\section*{Onde este resultado entra}

A pergunta ``quais moduladoras afetam quais coeficientes'' é a de identificar
os blocos $(\cv,\md)$ com $\gcomp{\cv}{\md} \equiv 0$. A decisão D28 fixou
duas frentes: a propriedade oráculo da variante em grupos, que pede uma
condição de irrepresentabilidade no desenho de produtos e é sondada em E1.7a,
e a \emph{limiarização}, que é este arquivo. As duas não competem; esta é a
que entra no artigo de qualquer forma, por três razões.

\begin{enumerate}[label=(\alph*), leftmargin=*, nosep]
  \item Ela consome o Corolário~4 de E1.6 \emph{como está} e não introduz
    hipótese de desenho nenhuma. A única hipótese nova é sobre o
    \emph{alvo} --- os blocos ativos têm norma acima da taxa de estimação ---
    e não sobre $(\bX,\bU)$.
  \item Ela vale para o \emph{LASSO puro}, que é o estimador base do projeto
    (D3) e o único que tem taxa provada aqui. A variante em grupos, que
    zera blocos sozinha, não tem teoria neste projeto
    (\texttt{docs/selecao-estrutura.md}).
  \item A estatística que ela limiariza já existe no código:
    $\norm{\thetahat_{\cv\md,\cdot}}_{2}$, que é o que
    \texttt{wafc\_blocks()} devolve, é \emph{exatamente}
    $\lVert\ghat{\cv}{\md}\rVert_{L_{2}[0,1]}$ pela ortonormalidade da base
    (Observação~\ref{rem:codigo}).
\end{enumerate}

O enunciado é deliberadamente modesto, e o arquivo o diz com todas as letras:
\emph{isto é estimação seguida de limiar, e não seleção pelo estimador}. Nada
aqui afirma que o LASSO escolhe os blocos certos; o que se prova é que, tendo
estimado bem cada componente, um corte na norma de bloco recupera a estrutura,
desde que a estrutura seja detectável na resolução em que se estimou. A
hipótese de separação aparece como hipótese, numerada, e a
Observação~\ref{rem:exata} mostra que a janela de limiares do
Lema~\ref{lem:sanduiche} é exata: quando a separação cai ao dobro do erro,
existe configuração em que \emph{nenhum} limiar acerta.

São dois enunciados. O Lema~\ref{lem:sanduiche} é determinístico e não sabe
o que é um LASSO: dado qualquer estimador cujo erro por bloco seja dominado
por $D$, ele sanduicha o conjunto limiarizado entre dois conjuntos de nível
da verdade. O Corolário~\ref{cor:limiar} instala nele a taxa do Corolário~4 e
produz $\Prob(\Shat = S) \to 1$.

\section{Objetos}

Mantém-se tudo de E1.6, \texttt{05-taxas.tex}, §1: o modelo, a amostra
i.i.d., a partição $\bZ = [\Amat\;\Bmat]$ na ordem de D12 com
$d = pq\NJ$, o estimador
\[
  \betahat = (\chat,\thetahat)
  \;\in\; \operatorname*{arg\,min}_{(\bc,\btheta)}
  \bigl\{\normn{\bY - \Amat\bc - \Bmat\btheta}^{2} + 2\pen\norm{\btheta}_{1}\bigr\},
\]
o resíduo determinístico $\bb = f - f_{J}$, a perfilagem
$\gtil = \lmin(\Btil\T\Btil/n)$, $\gamma = \kappa_{1}c_{U}$ e
$\mathcal{B}_{J}^{2} = (pq)^{2}C_{X}^{2}C_{U}C_{g}^{2}(1-2^{-2\seff})^{-1}
2^{-2J\seff}$. A penalidade é a do Corolário~2 de E1.5,
\[
  \pen_{n} \;=\; 2\sigma B_{X}\sqrt{2C_{U}}\;\sqrt{\frac{2L_{n}}{n}},
  \qquad L_{n} = \log(2d_{n}/\alpha),
\]
e a resolução é a do Teorema~2 de E1.6,
$J_{n} = \min\{J : 2^{J}\ge(n/\log n)^{1/(2\seff+1)}\}$.

Os objetos próprios deste arquivo são quatro. As componentes estimadas e as
suas normas,
\[
  \ghat{\cv}{\md} \;=\; \sum_{\lev<J_{n},\,\tsl}
  \widehat\theta_{\cv\md,\lev\tsl}\,\wav{\lev}{\tsl},
  \qquad
  \Nhat{\cv}{\md} \;=\; \bigl\lVert \ghat{\cv}{\md}\bigr\rVert_{L_{2}[0,1]},
  \qquad
  \Ntrue{\cv}{\md} \;=\; \bigl\lVert \gcomp{\cv}{\md}\bigr\rVert_{L_{2}[0,1]};
\]
a estrutura verdadeira e a estimada, esta indexada pelo limiar $t \ge 0$,
\[
  S \;=\; \bigl\{(\cv,\md) :\; \gcomp{\cv}{\md} \not\equiv 0
  \text{ em } L_{2}[0,1]\bigr\}
  \;=\; \bigl\{(\cv,\md) :\; \Ntrue{\cv}{\md} > 0 \bigr\},
  \qquad
  \Shat(t) \;=\; \bigl\{(\cv,\md) :\; \Nhat{\cv}{\md} > t \bigr\};
\]
o erro máximo por bloco e a separação,
\[
  \Delta \;=\; \max_{\cv,\md}
  \bigl\lVert \ghat{\cv}{\md} - \gcomp{\cv}{\md}\bigr\rVert_{L_{2}[0,1]},
  \qquad
  \sepmin \;=\; \min_{(\cv,\md)\in S} \Ntrue{\cv}{\md};
\]
e a taxa por componente do Corolário~4 de E1.6,
\begin{equation}
  \label{eq:rho}
  \rn_{n} \;=\; \Bigl(\frac{J_{n}}{n}\Bigr)^{\seff/(2\seff+1)}
  \;\asymp\; \Bigl(\frac{\log n}{n}\Bigr)^{\seff/(2\seff+1)} .
\end{equation}

Duas convenções. A primeira é que $S$ é um conjunto de \emph{blocos}, não de
coeficientes: $|S| \le pq$, e nada aqui diz respeito a quais
$\wav{\lev}{\tsl}$ dentro de um bloco são não nulos. A segunda é que a
desigualdade que define $\Shat(t)$ é estrita, o que só importa em empates de
medida nula e é a convenção que a conferência numérica implementa.

\section{Hipóteses}

Além das hipóteses de E1.3 a E1.6, que são citadas onde entram, este arquivo
usa uma, e ela é sobre o alvo, não sobre o desenho. O contador de hipóteses é
local ao arquivo, como em E1.5 e E1.6.

\begin{assumption}[separação em nível de bloco]\label{ass:sep}
Na sequência de problemas da Hipótese~1 de E1.6, o alvo pode variar com $n$;
escreva $S_{n}$ e $\sepmin_{n} = \min_{(\cv,\md)\in S_{n}}
\lVert g^{(n)}_{\cv\md}\rVert_{L_{2}[0,1]}$, com $\sepmin_{n} = +\infty$ se
$S_{n} = \emptyset$. Supõe-se
\[
  \frac{\sepmin_{n}}{\rn_{n}} \;\longrightarrow\; \infty ,
\]
com $\rn_{n}$ de \eqref{eq:rho}: \emph{o menor bloco ativo tem norma de ordem
estritamente maior que a taxa por componente}.
\end{assumption}

Três leituras, que a Observação~\ref{rem:fixa} retoma. (i) Sob alvo fixo, e
portanto sob a Hipótese~1 de E1.6 lida literalmente, $\sepmin_{n} = \sepmin$ é
constante positiva e a hipótese é automática, porque $\rn_{n}\to 0$. (ii) Ela
tem conteúdo em arranjo triangular, onde o menor bloco ativo pode encolher com
$n$; aí ela diz exatamente quão devagar. (iii) Ela é hipótese sobre
$\gcomp{\cv}{\md}$ e não sobre $(\bX,\bU)$: nenhuma condição de desenho nova
entra neste arquivo, e em particular nenhuma irrepresentabilidade.

\section{Os dois enunciados}

\begin{lemma}[limiarização de um estimador com erro uniforme sobre blocos]
\label{lem:sanduiche}
Sejam $\{\gcomp{\cv}{\md}\}$ e $\{\ghat{\cv}{\md}\}$ famílias quaisquer de
funções de $L_{2}[0,1]$ indexadas por $(\cv,\md)$, com $\Nhat{\cv}{\md}$,
$\Ntrue{\cv}{\md}$, $S$, $\Shat(t)$, $\Delta$ e $\sepmin$ definidos a partir
delas como na §1, e seja $D \ge \Delta$. Então, para todo $t \ge 0$,
\begin{equation}
  \label{eq:sanduiche}
  \bigl\{(\cv,\md):\; \Ntrue{\cv}{\md} > t + D\bigr\}
  \;\subseteq\; \Shat(t) \;\subseteq\;
  \bigl\{(\cv,\md):\; \Ntrue{\cv}{\md} > t - D\bigr\} .
\end{equation}
Em particular:
\begin{enumerate}[label=\textup{(\roman*)}, leftmargin=*, nosep]
  \item \textup{(sem falso positivo; não usa separação)} se $t \ge D$, então
    $\Shat(t) \subseteq S$;
  \item \textup{(sem falso negativo; usa separação)} se $t < \sepmin - D$,
    então $S \subseteq \Shat(t)$;
  \item \textup{(acerto exato)} se $D \le t < \sepmin - D$, então
    $\Shat(t) = S$. Tal $t$ existe se e só se $\sepmin > 2D$.
\end{enumerate}
\end{lemma}

O lema não usa nada do estimador além da cota $\Delta \le D$: vale para o
LASSO, para a variante em grupos, para um estimador por splines, e é por isso
que a única coisa que E1.6 precisa fornecer é a taxa por componente.

\begin{corollary}[seleção de estrutura por limiarização]\label{cor:limiar}
Nas condições do Teorema~2 de E1.6 (hipóteses de E1.3, E1.4 e E1.5, a
Hipótese~1 de E1.6, o $\pen_{n}$ e o
$J_{n}$ daquele arquivo):
\begin{enumerate}[label=\textup{(\roman*)}, leftmargin=*]
  \item \textup{(não assintótico)} Sejam $\alpha,\alpha'\in(0,1)$ e
    \begin{equation}
      \label{eq:Dn}
      D_{n}^{2}
      \;=\; \frac{512\,\pen_{n}^{2}\,d_{n}}{\gamma^{2}}
      \;+\; \frac{64\,\mathcal{B}_{J_{n}}^{2}}{\alpha'\gamma}
      \;+\; \frac{2\,C_{g}^{2}}{1-2^{-2\seff}}\,2^{-2J_{n}\seff} .
    \end{equation}
    Então, num evento de probabilidade ao menos
    $1-\alpha-\alpha'-\Prob\{\lmin(\Gramhat)<\gamma/2\}
    -\Prob\{\smax^{2}>2B_{X}^{2}C_{U}\}$, vale $\Delta \le D_{n}$ e portanto,
    pelo Lema~\ref{lem:sanduiche}(iii),
    \[
      \Shat(t) = S \qquad\text{para todo } t \text{ com }
      D_{n} \le t < \sepmin - D_{n} ,
    \]
    janela não vazia sempre que $\sepmin > 2D_{n}$. As duas últimas
    probabilidades tendem a zero pela Proposição~3(iv) de E1.4 e pelo Lema~6
    de E1.5.
  \item \textup{(triagem, sem hipótese de separação)} $\Delta = O_{p}(\rn_{n})$
    e, para qualquer $t_{n}$ com $t_{n}/\rn_{n} \to \infty$,
    \[
      \Prob\bigl\{\Shat(t_{n}) \subseteq S\bigr\} \;\longrightarrow\; 1 .
    \]
    \emph{Nenhuma hipótese sobre $\sepmin_{n}$ é usada aqui}: eliminar os
    blocos nulos é de graça.
  \item \textup{(acerto exato)} Sob a Hipótese~\ref{ass:sep}, e para qualquer
    $t_{n}$ com
    \[
      \frac{t_{n}}{\rn_{n}} \to \infty
      \qquad\text{e}\qquad
      \limsup_{n} \frac{t_{n}}{\sepmin_{n}} < 1 ,
    \]
    vale $\Prob\{\Shat(t_{n}) = S\} \to 1$. Tal $t_{n}$ existe sempre sob a
    Hipótese~\ref{ass:sep}, e a escolha canônica é a média geométrica
    $t_{n} = \sqrt{\rn_{n}\sepmin_{n}}$.
\end{enumerate}
\end{corollary}

\section{Provas}

\subsection*{Lema~\ref{lem:sanduiche}}

Pela desigualdade triangular em $L_{2}[0,1]$,
\begin{equation}
  \label{eq:triangular}
  \bigl|\Nhat{\cv}{\md} - \Ntrue{\cv}{\md}\bigr|
  \;\le\; \bigl\lVert\ghat{\cv}{\md} - \gcomp{\cv}{\md}\bigr\rVert_{L_{2}[0,1]}
  \;\le\; \Delta \;\le\; D
  \qquad\text{para todo }(\cv,\md).
\end{equation}
Inclusão da esquerda de \eqref{eq:sanduiche}: se $\Ntrue{\cv}{\md} > t+D$,
então $\Nhat{\cv}{\md} \ge \Ntrue{\cv}{\md} - D > t$, logo
$(\cv,\md)\in\Shat(t)$. Inclusão da direita: se $(\cv,\md)\in\Shat(t)$, então
$\Ntrue{\cv}{\md} \ge \Nhat{\cv}{\md} - D > t - D$.

Para (i), tome $t \ge D$: a inclusão da direita dá
$\Shat(t) \subseteq \{\Ntrue{\cv}{\md} > t - D\} \subseteq
\{\Ntrue{\cv}{\md} > 0\} = S$, porque $t - D \ge 0$. Para (ii), tome
$t < \sepmin - D$: todo $(\cv,\md)\in S$ tem
$\Ntrue{\cv}{\md} \ge \sepmin > t + D$, e a inclusão da esquerda dá
$(\cv,\md)\in\Shat(t)$. (iii) é a conjunção das duas, e o intervalo
$[D,\sepmin-D)$ é não vazio se e só se $D < \sepmin - D$. \qed

\subsection*{Corolário~\ref{cor:limiar}}

\emph{(i).} Intersecte os quatro eventos da prova do Corolário~1 de E1.5:
$\mathcal{T}$, de probabilidade ao menos $1-\alpha$;
$\{\smax^{2}\le 2B_{X}^{2}C_{U}\}$, que torna $\pen_{n}$ admissível, isto é
$\pen_{n}\ge 2\pen_{0}$; $\{\lmin(\Gramhat)\ge\gamma/2\}$, que com o
complemento de Schur do Lema~4 de E1.5 dá $\gtil \ge \gamma/2$; e
$\{\normn{\bb}^{2} \le \mathcal{B}_{J_{n}}^{2}/\alpha'\}$, que é Markov
aplicado a $\E\normn{\bb}^{2} \le \mathcal{B}_{J_{n}}^{2}$ (Corolário~1 de
E1.3). Nesse evento, o Teorema~1(ii) de E1.5 com $s_{0} \le d_{n}$ dá
\[
  \norm{v}_{2}^{2}
  \;\le\; \frac{64\pen_{n}^{2}d_{n}}{\gtil^{2}} + \frac{16\normn{\bb}^{2}}{\gtil}
  \;\le\; \frac{256\pen_{n}^{2}d_{n}}{\gamma^{2}}
  + \frac{32\,\mathcal{B}_{J_{n}}^{2}}{\alpha'\gamma},
  \qquad v = \thetahat - \thetastar .
\]
Fixe agora um bloco. Pela ortonormalidade das $\{\wav{\lev}{\tsl}\}_{\lev<J_{n}}$
em $L_{2}[0,1]$,
$\lVert\ghat{\cv}{\md} - \PiJ\gcomp{\cv}{\md}\rVert_{L_{2}[0,1]}
= \norm{v_{\cv\md,\cdot}}_{2} \le \norm{v}_{2}$, e o Lema~1 de E1.3 dá
$\lVert\PiJ\gcomp{\cv}{\md} - \gcomp{\cv}{\md}\rVert^{2}_{L_{2}[0,1]}
\le C_{g}^{2}(1-2^{-2\seff})^{-1}2^{-2J_{n}\seff}$. Por
$(x+y)^{2}\le 2x^{2}+2y^{2}$,
\[
  \Delta^{2} \;\le\; 2\norm{v}_{2}^{2}
  + \frac{2C_{g}^{2}}{1-2^{-2\seff}}2^{-2J_{n}\seff}
  \;\le\; D_{n}^{2},
\]
que é \eqref{eq:Dn}. O Lema~\ref{lem:sanduiche}(iii) com $D = D_{n}$ fecha.

\emph{(ii).} Na resolução $J_{n}$ do Teorema~2 de E1.6, as três parcelas de
\eqref{eq:Dn} são da mesma ordem $\rn_{n}^{2}$: a primeira porque
$\pen_{n}^{2}d_{n} \asymp (J_{n}/n)\,pq2^{J_{n}} \asymp (J_{n}/n)^{2\seff/(2\seff+1)}$
pelo Teorema~2(iii) de E1.6, e a segunda e a terceira porque
$2^{-2J_{n}\seff} \asymp (\log n/n)^{2\seff/(2\seff+1)}$ pela definição de
$J_{n}$. Logo $D_{n} \le c(\alpha')\rn_{n}$ com $c(\alpha')<\infty$. Dado
$\epsilon>0$, tome $\alpha = \alpha' = \epsilon/4$ e $n$ grande o bastante
para que as duas probabilidades residuais de (i) somem menos de
$\epsilon/2$; então $\Prob(\Delta > c(\epsilon/4)\rn_{n}) < \epsilon$, que é
$\Delta = O_{p}(\rn_{n})$. Se $t_{n}/\rn_{n}\to\infty$, então
$\Prob(t_{n} \ge \Delta) \to 1$, e nesse evento o
Lema~\ref{lem:sanduiche}(i) dá $\Shat(t_{n})\subseteq S$.

\emph{(iii).} Sob a Hipótese~\ref{ass:sep} e com
$\limsup t_{n}/\sepmin_{n} < 1$, existem $\eta>0$ e $n_{0}$ tais que
$\sepmin_{n} - t_{n} \ge \eta\,\sepmin_{n}$ para $n \ge n_{0}$, e portanto
$(\sepmin_{n}-t_{n})/\rn_{n} \ge \eta\,\sepmin_{n}/\rn_{n} \to\infty$.
Assim $\Prob(\Delta \le t_{n}) \to 1$ e
$\Prob(\Delta < \sepmin_{n}-t_{n}) \to 1$; na interseção dos dois eventos,
o Lema~\ref{lem:sanduiche}(iii) com $D = \Delta$ dá $\Shat(t_{n}) = S$.
Quanto à existência: com $t_{n} = \sqrt{\rn_{n}\sepmin_{n}}$ tem-se
$t_{n}/\rn_{n} = \sqrt{\sepmin_{n}/\rn_{n}} \to\infty$ e
$t_{n}/\sepmin_{n} = \sqrt{\rn_{n}/\sepmin_{n}} \to 0$. \qed

\section{Observações}

\begin{remark}[a janela do Lema~\ref{lem:sanduiche} é exata]\label{rem:exata}
A condição $\sepmin > 2D$ não é artefato da prova. Tome dois blocos, um ativo
com $\Ntrue{1}{1} = 2D$ e um nulo com $\Ntrue{2}{1} = 0$, e estimativas
$\Nhat{1}{1} = D$ e $\Nhat{2}{1} = D$, compatíveis com $\Delta \le D$. As
duas estatísticas empatam, e nenhum $t$ separa. Isto diz que o \emph{argumento}
não se estende abaixo de $\sepmin = 2D$; não diz que nenhum procedimento
consegue, o que seria uma cota inferior e não é feito aqui nem em E1.6
(``o que isto não cobre'' daquele arquivo).
\end{remark}

\begin{remark}[a hipótese de separação sob alvo fixo]\label{rem:fixa}
A Hipótese~1 de E1.6 fixa $p$, $q$, $s$, $\pi$ e $C_{g}$, e se o alvo também
for fixo então $S$ e $\sepmin$ não dependem de $n$; como $\rn_{n}\to 0$, a
Hipótese~\ref{ass:sep} é automática e o
Corolário~\ref{cor:limiar}(iii) vale para \emph{qualquer} $t_{n}\to 0$ com
$t_{n}/\rn_{n}\to\infty$. É a leitura que a conferência numérica realiza, e é
o enunciado que um leitor de simulação espera. A forma em arranjo triangular é
a que carrega informação: ela diz que a fronteira da detecção é a taxa de
estimação por componente, e é ela que torna a hipótese uma hipótese, e não uma
formalidade.
\end{remark}

\begin{remark}[é estimação seguida de limiar, e não seleção pelo estimador]
\label{rem:nao-selecao}
Nada aqui afirma que o LASSO seleciona. O suporte de $\thetahat$ não tem
relação com $S$: sob a Hipótese de Besov de E1.3 o oráculo do espaço de
aproximação é genericamente denso (E1.6, ``onde este resultado entra''), e na
prática o LASSO deixa coeficientes não nulos em blocos nulos --- a conferência de E2.2
mediu, com cerca de $99$ não nulos, \emph{zero} dos três blocos nulos zerados
pelo LASSO contra três pelo sparse group LASSO. O que o
Corolário~\ref{cor:limiar} diz é outra coisa: a \emph{norma} do bloco, e não o
seu suporte, é pequena nos blocos nulos, e um corte na norma recupera $S$. O
procedimento é, portanto, de dois passos, e o enunciado do artigo tem de
dizê-lo.
\end{remark}

\begin{remark}[a estatística é a que o código já devolve]\label{rem:codigo}
Como $\{\wav{\lev}{\tsl}\}_{\lev<J_{n}}$ é ortonormal em $L_{2}[0,1]$,
\[
  \Nhat{\cv}{\md}
  \;=\; \Bigl\lVert \sum_{\lev<J_{n},\tsl}\widehat\theta_{\cv\md,\lev\tsl}
  \wav{\lev}{\tsl}\Bigr\rVert_{L_{2}[0,1]}
  \;=\; \norm{\thetahat_{\cv\md,\cdot}}_{2},
\]
que é a coluna \texttt{norm} de \texttt{wafc\_blocks()} (E2.2). A identidade é
exata no regime da teoria, $\eps = 0$ (D26), e a conferência a verifica com
erro relativo máximo de $1{,}4\times10^{-9}$ nos seis blocos. Com a margem do código
(\texttt{rescale = TRUE}, $\eps$ hoje em $0{,}05$), a base é ortonormal na
variável reescalada, de modo que $\norm{\thetahat_{\cv\md,\cdot}}_{2}$ é a
norma de $\ghat{\cv}{\md}$ no intervalo \emph{alargado} dividida por
$\sqrt{\text{scale}}$. O fator é comum a todos os blocos que compartilham
$U_{\md}$ e portanto não muda a ordenação das estatísticas dentro de uma
coluna $\md$; entre colunas com amplitudes amostrais diferentes, muda, e é
ponto para E2.5 calibrar.
\end{remark}

\begin{remark}[o que muda se a variante em grupos ganhar taxa]\label{rem:grupos}
O Lema~\ref{lem:sanduiche} é determinístico e cego ao estimador. Se E1.7a ou
um trabalho posterior der uma taxa por componente para o sparse group LASSO,
o Corolário~\ref{cor:limiar} vale para ele trocando $\rn_{n}$ pela taxa nova,
sem reescrever uma linha da prova. O que \emph{não} se obtém assim é a
propriedade oráculo de seleção da saída (a): esta continua sendo uma
afirmação sobre um procedimento de dois passos, e aquela é sobre o estimador
em um passo.
\end{remark}

\section{O que foi transposto e o que é novo}

\begin{itemize}[nosep, leftmargin=*]
  \item \textbf{Transposto:} a ideia de ``estimar bem e limiarizar seleciona''
    é clássica em LASSO, e a referência mais próxima em forma é van de Geer,
    Bühlmann e Zhou (2011), que estudam o LASSO limiarizado exatamente como um
    segundo passo sobre um estimador com taxa; em modelos aditivos, Huang,
    Horowitz e Wei (2010) usam uma condição de separação em nível de
    \emph{componente} no mesmo papel que a Hipótese~\ref{ass:sep} aqui.
  \item \textbf{Novo, no sentido em que este projeto usa a palavra:} o
    enunciado no \emph{desenho de produtos} $X_{\cv}\wav{\lev}{\tsl}(U_{\md})$
    com $q\ge2$ moduladoras, em que a unidade limiarizada é o bloco
    $(\cv,\md)$ --- isto é, a resposta à pergunta ``$U_{\md}$ modula
    $\fcoef{\cv}$?'' --- e não uma variável nem uma componente aditiva
    isolada. E o fato de valer para o \emph{LASSO puro}, sem
    irrepresentabilidade e sem condição de desenho além das de E1.4, que já
    estão pagas.
  \item \textbf{Novo aqui:} a separação entre as duas metades do problema no
    Corolário~\ref{cor:limiar}(ii) e (iii). Eliminar os blocos nulos não
    custa hipótese nenhuma; toda a hipótese de separação é gasta em não
    perder blocos ativos pequenos. A conferência mede as duas curvas
    separadamente, e é essa assimetria que explica por que o acerto observado
    é alto mesmo quando a janela do Lema~\ref{lem:sanduiche} é vazia.
\end{itemize}

\section{O que isto não cobre}

\begin{itemize}[nosep, leftmargin=*]
  \item \textbf{$t_{n}$ depende de constantes desconhecidas.} A janela de (i)
    é escrita em $D_{n}$, que carrega $\sigma$, $\gamma = \kappa_{1}c_{U}$,
    $C_{g}$, $\seff$ e $\mathcal{B}_{J_{n}}$; e a de (iii), em $\rn_{n}$ e
    $\sepmin_{n}$. Nada aqui diz como escolher $t_{n}$ com os dados. A
    escolha empírica é de E2.4 e E2.5, e a conferência desta etapa entrega
    apenas a curva de acerto contra o limiar e uma regra invariante de escala
    para servir de ponto de partida.
  \item \textbf{Não é cota inferior.} A Observação~\ref{rem:exata} mostra que
    o argumento para em $\sepmin = 2D$; não se prova que a separação seja
    necessária para algum procedimento. A cota inferior para $q\ge2$ não
    existe neste projeto (E1.9).
  \item \textbf{Não é a propriedade oráculo da saída (a).} O estimador não
    seleciona; o procedimento de dois passos seleciona
    (Observação~\ref{rem:nao-selecao}). Se E1.7a fechar, os dois enunciados
    coexistem, como D28 previu.
  \item \textbf{Nada sobre o suporte dentro do bloco.} Quais
    $\wav{\lev}{\tsl}$ são não nulos, e se o conjunto deles tem algum
    significado, fica fora. A dimensão efetiva do Corolário~5 de E1.6 é uma
    afirmação de taxa, não de recuperação de suporte, e continua sendo.
  \item \textbf{O alvo limiarizado é $\PiJ\gcomp{\cv}{\md}$, não
    $\gcomp{\cv}{\md}$.} A estatística vive na resolução $J_{n}$, e o viés de
    aproximação entra em $D_{n}$ pela terceira parcela de \eqref{eq:Dn}. Nas
    resoluções que a regra $J_{n}$ prescreve para os $n$ do piloto essa
    parcela \emph{domina} as outras duas, o que a conferência mede: a janela
    do Lema~\ref{lem:sanduiche} chega a ser vazia enquanto o acerto observado
    é 1, porque o viés encolhe todas as normas na mesma direção e não cria
    falso positivo.
  \item \textbf{$\sigma$ conhecido, $\pen$ determinístico, $\eps = 0$,
    $p$ e $q$ fixos}, pelas mesmas razões de E1.5 e E1.6. Com $pq$ crescendo,
    $\Delta \le \norm{v}_{2}$ continua valendo sem união sobre blocos, mas as
    constantes de E1.3 e E1.4 carregam $pq$ e a janela estreita.
  \item \textbf{A variante em grupos continua sem teoria} neste projeto
    (Observação~\ref{rem:grupos}); a comparação numérica da conferência é
    empírica e não sustenta enunciado.
\end{itemize}

\section{Conferência numérica}

\texttt{check/06-selecao-limiar.R} (imprime \texttt{OK}; cerca de $7$~min
com o padrão de $20$ réplicas). Cenários, desenho e estimador são os de
\texttt{wafc/R/}, com $p = 3$, $q = 2$, Daublets de filtro $8$, \texttt{snr}
$= 4$, \texttt{rescale = FALSE} (o regime $\eps = 0$ da teoria, D26) e a base
avaliada por tabela fixa nas partes de simulação (D31), com avaliação exata
na parte que mede a identidade de normas. As cinco partes e o que saiu em
2026-09-20 estão em \texttt{docs/handoff-E1.7c.md}; os números que importam:

\begin{itemize}[nosep, leftmargin=*]
  \item \textbf{Parte I, o Lema~\ref{lem:sanduiche}.} O sanduíche
    \eqref{eq:sanduiche} e a janela de (iii) valem em $2000$ configurações
    aleatórias de até $12$ blocos, com $21$ limiares em cada uma: zero
    falhas nas duas afirmações. A configuração da Observação~\ref{rem:exata}
    é construída e verificada: com $\sepmin = 2D$ as duas estatísticas
    empatam em $D$ e nenhum dos $201$ limiares testados acerta.
  \item \textbf{Parte II, a estatística.} Com \texttt{rescale = FALSE} e
    avaliação exata da base, $\Nhat{\cv}{\md}$ calculada por quadratura e
    $\norm{\thetahat_{\cv\md,\cdot}}_{2}$ coincidem nos seis blocos com erro
    relativo máximo de $1{,}4\times10^{-9}$. Com a margem do código
    ($\eps = 0{,}05$, \texttt{scale} $= 1{,}108$), a norma no intervalo
    \emph{observado} é de $0{,}938$ a $0{,}998$ vezes a estatística do
    código: a diferença é a energia que $\ghat{\cv}{\md}$ põe na margem.
  \item \textbf{Parte III, a decomposição de $\Delta_{n}$.} Com
    \texttt{cv.min}, a parcela de estimação vai de $0{,}184$ a $0{,}076$ no
    cenário suave e de $0{,}409$ a $0{,}213$ no não homogêneo quando $n$ vai
    de $250$ a $2000$, com inclinações log-log contra
    $(\log n/n)^{\seff/(2\seff+1)}$ de $1{,}38$ e $1{,}56$: o erro realizado
    cai \emph{mais rápido} que a cota, como em E1.6. A parcela de viés, em
    $n = 2000$, é $0{,}033$ no suave e $0{,}635$ no não homogêneo, onde
    \emph{domina} $\Delta_{n}$.
  \item \textbf{Parte III, a janela e o acerto.} $\sepmin = 1$ nos dois
    cenários por construção, mas a separação efetiva
    $\min_{S}\lVert\PiJ\gcomp{\cv}{\md}\rVert$ é $0{,}999$ no suave e
    $0{,}778$ no não homogêneo em $n = 2000$. A janela $\sepmin > 2\Delta_{n}$
    existe em todas as células do cenário suave e em \emph{nenhuma} do não
    homogêneo; ainda assim, $\Prob(\Shat = S)$ no melhor limiar da grade é
    $1{,}00$ em todas as células com \texttt{cv.min} a partir de $n = 500$, e
    a fração de réplicas em que \emph{algum} limiar acerta é $1{,}00$ (a
    única exceção é $0{,}95$, em $n = 250$ no não homogêneo). É a assimetria
    do Corolário~\ref{cor:limiar}(ii) contra (iii): o viés de aproximação
    encolhe todas as normas na mesma direção e não cria falso positivo.
  \item \textbf{Parte III, a curva.} Em $n = 2000$ o platô de acerto é largo:
    $\Prob(\Shat(t) = S) = 1$ para $t \in [0{,}062;\,0{,}590]$ no cenário
    suave e $t \in [0{,}127;\,0{,}590]$ no não homogêneo. À esquerda do platô
    entram os três blocos nulos, à direita saem os três ativos.
  \item \textbf{Parte III, a regra da teoria.} O par $(J_{n},\pen_{n})$ do
    Teorema~2 e do Corolário~2, com $\sigma$ conhecido, \emph{não} serve para
    limiarizar nos $n$ do piloto: ele zera quase tudo, o melhor limiar é o
    piso da grade ($0{,}002$), e em $n = 250$ no cenário não homogêneo o
    acerto é $0$. É a mesma leitura pré-assintótica que E1.6 registrou em
    taxa e E2.3 em sintonia, agora em seleção.
  \item \textbf{Parte IV, a Hipótese~S medida.} Em $J = 4$ e $n = 1000$, onde
    a parcela de viés é pequena, $\Delta_{n} \approx 0{,}10$ e quase não
    depende de $\sepmin$. Fazendo $\sepmin$ percorrer $1$, $1/2$, $1/4$ e
    $1/8$, a razão $\sepmin/\Delta_{n}$ cai de $9{,}00$ a $1{,}30$, a janela
    $\sepmin > 2\Delta_{n}$ existe nos três primeiros e some no quarto, e o
    acerto cai de $1{,}00$ a $0{,}85$. É a hipótese de separação vista
    fechar.
  \item \textbf{Parte V, contra a variante em grupos.} Em $n = 1000$,
    $J = 4$, $10$ réplicas e as \emph{mesmas dobras}: o LASSO limiarizado
    acerta a estrutura em $1{,}00$ das réplicas nos dois cenários com o
    melhor limiar da grade, e em $1{,}00$ e $0{,}90$ com a regra invariante
    de escala $t = 0{,}15\max_{\cv\md}\Nhat{\cv}{\md}$, que não é teoria e
    serve só de ponto de partida para E2.4. O sparse group LASSO em
    \texttt{lambda.min}, lido pelos blocos que ele zera, acerta em
    \emph{zero} réplicas nos dois cenários. O RMSE fora da amostra é
    $0{,}770$ contra $0{,}776$ no suave e $1{,}310$ contra $1{,}298$ no não
    homogêneo, com $\sigma \approx 0{,}76$: a estrutura sai de graça em
    predição. A comparação é empírica e não sustenta enunciado
    (Observação~\ref{rem:grupos}); o que ela mostra é que a pergunta de E2.5
    deixou de ser ``grupos ou nada''.
\end{itemize}

\section*{Referências}
\sloppy

Referências fora de \texttt{docs/referencias-verificadas.bib} vão ao handoff.

\begin{itemize}[nosep, leftmargin=*]
  \item Bühlmann, P.\ e van de Geer, S. (2011). \emph{Statistics for
    High-Dimensional Data}. Springer. (\texttt{buhlmann2011statistics},
    verificada em L1.) Fonte do Teorema~1 de E1.5, que é tudo o que este
    arquivo consome de LASSO. O tratamento do LASSO limiarizado está no
    livro. \textsc{[verificar numeração da seção]}
  \item van de Geer, S., Bühlmann, P.\ e Zhou, S. (2011). The adaptive and
    the thresholded Lasso for potentially misspecified models (and a lower
    bound for the Lasso). \emph{Electronic Journal of Statistics} 5.
    \texttt{10.1214/11-EJS624}. Conferida no Crossref em 2026-09-20 (autores,
    título, veículo, volume, ano e DOI); não está no \texttt{.bib} e o
    intervalo de páginas não consta do registro
    \textsc{[verificar páginas]}. Referência de forma do procedimento de
    dois passos.
  \item Huang, J., Horowitz, J.~L.\ e Wei, F. (2010). Variable selection in
    nonparametric additive models. \emph{The Annals of Statistics} 38(4),
    2282--2313. (\texttt{Huang-Horowitz-Wei-2010}, verificada em L1.)
    Condição de separação em nível de componente no mesmo papel da
    Hipótese~\ref{ass:sep}. \textsc{[verificar numeração do resultado]}
  \item Wei, F., Huang, J.\ e Li, H. (2011). Variable selection and
    estimation in high-dimensional varying-coefficient models.
    \emph{Statistica Sinica} 21(4), 1515--1540.
    (\texttt{Wei-Huang-Li-2011}, verificada em L1.) A rota da saída (a), que
    este arquivo não segue.
  \item Klopp, O.\ e Pensky, M. (2015). Sparse high-dimensional varying
    coefficient model: nonasymptotic minimax study. \emph{The Annals of
    Statistics} 43(3), 1273--1299. (\texttt{Klopp-Pensky-2015}, verificada em
    L1.) O vizinho de E1.6; não trata de seleção de estrutura.
\end{itemize}

\end{document}
